\subsubsection{Generic Controlled Gate (CU)}\label{sec:generic-controlled-gate-cu}

The \definition{Generic Controlled Gate (CU)} is a \textbf{two-qubit gate} composed of:
\begin{itemize}
    \item A \textbf{control qubit} that determines \hl{whether} the operation is \hl{applied or not} (first qubit).
    \item A \textbf{target qubit} \hl{on which the operation is applied} if the control qubit is in the state $\ket{1}$ (second qubit).
    \item A \textbf{unitary operation} $U$ that is \hl{applied to the target qubit} when the control qubit is in the state $\ket{1}$.
\end{itemize}
It \textbf{generalizes the CNOT gate} (\autopageref{sec:controlled-not-cnot-gate}, $U = X$) by allowing \textbf{any unitary operation} $U$ to be applied to the target qubit, rather than just the NOT operation.

\highspace
\begin{flushleft}
    \textcolor{Green3}{\faIcon{book} \textbf{Action on Basis States}}
\end{flushleft}
The action of the CU gate on the computational basis states can be described as follows:

\begin{table}[!htp]
    \centering
    \begin{tabular}{@{} c | c @{}}
        \toprule
        \textbf{Input} & \textbf{Output} \\
        \midrule
        $\ket{00}$ & $\ket{00}$ \\
        $\ket{01}$ & $\ket{01}$ \\
        $\ket{10}$ & $\ket{1} \otimes U\ket{0}$ \\
        $\ket{11}$ & $\ket{1} \otimes U\ket{1}$ \\
        \bottomrule
    \end{tabular}
\end{table}

\noindent
If the control qubit is $\ket{0}$, the target qubit remains unchanged. If the control qubit is $\ket{1}$, the unitary operation $U$ is applied to the target qubit, transforming $\ket{0}$ to $U\ket{0}$ and $\ket{1}$ to $U\ket{1}$.

\highspace
\begin{flushleft}
    \textcolor{Green3}{\faIcon{square-root-alt} \textbf{Matrix Representation}}
\end{flushleft}
The CU gate can be represented as a $4 \times 4$ unitary matrix in the computational basis $\{\ket{00}, \ket{01}, \ket{10}, \ket{11}\}$:
\begin{equation}\label{eq:generic-controlled-gate-matrix}
    \text{CU} = C_{U} =
    \begin{bmatrix}
        I & 0 \\
        0 & U
    \end{bmatrix} =
    \begin{bmatrix}
        1 & 0 & 0 & 0 \\
        0 & 1 & 0 & 0 \\
        0 & 0 & U_{00} & U_{01} \\
        0 & 0 & U_{10} & U_{11}
    \end{bmatrix}
\end{equation}
Where $U_{00}, U_{01}, U_{10}, U_{11}$ are the elements of the $2 \times 2$ unitary matrix $U$ that acts on the target qubit when the control qubit is $\ket{1}$. For the first two basis states $\ket{00}$ and $\ket{01}$, the identity operation $I$ is applied, leaving them unchanged. For the last two basis states $\ket{10}$ and $\ket{11}$, the unitary operation $U$ is applied to the target qubit, resulting in the transformation defined by $U$.

\newpage

\begin{flushleft}
    \textcolor{Green3}{\faIcon{book} \textbf{Action on General States}}
\end{flushleft}
Let's consider two general single-qubit states:
\begin{equation*}
    \ket{v_A} = a_0\ket{0} + a_1\ket{1} \quad \ket{v_B} = b_0\ket{0} + b_1\ket{1}
\end{equation*}
Combining these states into a two-qubit state, we have:
\begin{equation*}
    \ket{v} = \ket{v_A} \otimes \ket{v_B} = a_0\ket{0}\ket{v_B} + a_1\ket{1}\ket{v_B}
\end{equation*}
Applying the CU gate to this state, we get:
\begin{align*}
    \text{CU}\ket{v} &= \text{CU}(a_0\ket{0}\ket{v_B} + a_1\ket{1}\ket{v_B}) \\
    &= a_0\text{CU}(\ket{0}\ket{v_B}) + a_1\text{CU}(\ket{1}\ket{v_B}) \\
    &= \underbrace{a_0\ket{0}\ket{v_B}}_{\text{unchanged}} + \underbrace{a_1\ket{1}U\ket{v_B}}_{\text{transformed}}
\end{align*}
This shows that the CU gate applies the unitary operation $U$ to the target qubit only when the control qubit is in the state $\ket{1}$, while leaving the target qubit unchanged when the control qubit is in the state $\ket{0}$. Each component of the superposition evolves independently according to the value of the control qubit, resulting in a new superposition where different branches have undergone different transformations.

\highspace
\begin{flushleft}
    \textcolor{Green3}{\faIcon{tools} \textbf{Circuit Representation}}
\end{flushleft}
In quantum circuit diagrams, the CU gate is typically represented with a \textbf{control dot on the control qubit} and a \textbf{box labeled with the unitary operation $U$ on the target qubit}. The control dot indicates that the operation is conditional on the state of the control qubit. The circuit representation of the CU gate is as follows:

\begin{center}
    \begin{quantikz}
        \lstick{$\ket{v_A}$} & \ctrl{1} & \qw \\
        \lstick{$\ket{v_B}$} & \gate{U} & \qw
    \end{quantikz}
\end{center}

\noindent
In summary, \hl{if the control qubit is $\ket{1}$, we apply $U$ to the target qubit; otherwise do nothing.}

\highspace
\begin{flushleft}
    \textcolor{Green3}{\faIcon{cogs} \textbf{Properties}}
\end{flushleft}
The CU gate has the same properties as the CNOT gate because it is a generalization of the CNOT gate. Specifically:
\begin{itemize}
    \item \important{Unitary}
    \begin{equation*}
        \text{CU}^\dagger \text{CU} = I
    \end{equation*}
    \item \important{Reversible}
    \begin{equation*}
        \text{CU}^{-1} = \text{CU}^\dagger
    \end{equation*}
    \item \important{Conditional Operation}
    \begin{equation*}
        \text{CU}\left(\ket{v_A} \otimes \ket{v_B}\right) \ne \text{CU}\ket{v_A} \otimes \text{CU}\ket{v_B}
    \end{equation*}
    It does not act independently on each qubit, but rather applies a transformation to the target qubit based on the state of the control qubit.
    \item \important{Entangling}. The CU gate can create entanglement between the control and target qubits if the unitary operation $U$ and the input state are chosen appropriately. For example, when $U = X$ (i.e., the CNOT gate), we have:
    \begin{equation*}
        \text{CNOT}(\ket{+}\ket{0}) = \frac{1}{\sqrt{2}}(\ket{00} + \ket{11})
    \end{equation*}
\end{itemize}